
\documentclass[a4paper,11pt]{article}
\usepackage[a4paper, margin=8em]{geometry}

% usa i pacchetti per la scrittura in italiano
\usepackage[french,italian]{babel}
\usepackage[T1]{fontenc}
\usepackage[utf8]{inputenc}
\frenchspacing 

% usa i pacchetti per la formattazione matematica
\usepackage{amsmath, amssymb, amsthm, amsfonts}

% usa altri pacchetti
\usepackage{gensymb}
\usepackage{hyperref}
\usepackage{standalone}

% imposta il titolo
\title{Appunti Comunicazioni Numeriche}
\author{Luca Seggiani}
\date{2026}

% imposta lo stile
% usa helvetica
\usepackage[scaled]{helvet}
% usa palatino
\usepackage{palatino}
% usa un font monospazio guardabile
\usepackage{lmodern}

\renewcommand{\rmdefault}{ppl}
\renewcommand{\sfdefault}{phv}
\renewcommand{\ttdefault}{lmtt}

% disponi il titolo
\makeatletter
\renewcommand{\maketitle} {
	\begin{center} 
		\begin{minipage}[t]{.8\textwidth}
			\textsf{\huge\bfseries \@title} 
		\end{minipage}%
		\begin{minipage}[t]{.2\textwidth}
			\raggedleft \vspace{-1.65em}
			\textsf{\small \@author} \vfill
			\textsf{\small \@date}
		\end{minipage}
		\par
	\end{center}

	\thispagestyle{empty}
	\pagestyle{fancy}
}
\makeatother

% disponi teoremi
\usepackage{tcolorbox}
\newtcolorbox[auto counter, number within=section]{theorem}[2][]{%
	colback=blue!10, 
	colframe=blue!40!black, 
	sharp corners=northwest,
	fonttitle=\sffamily\bfseries, 
	title=~\thetcbcounter: #2, 
	#1
}

% disponi definizioni
\newtcolorbox[auto counter, number within=section]{definition}[2][]{%
	colback=red!10,
	colframe=red!40!black,
	sharp corners=northwest,
	fonttitle=\sffamily\bfseries,
	title=~\thetcbcounter: #2,
	#1
}

% U.D.M
\newcommand{\amp}{\ensuremath{\mathrm{A}}}
\newcommand{\volt}{\ensuremath{\mathrm{V}}}
\newcommand{\meter}{\ensuremath{\mathrm{m}}}
\newcommand{\second}{\ensuremath{\mathrm{s}}}
\newcommand{\farad}{\ensuremath{\mathrm{F}}}
\newcommand{\henry}{\ensuremath{\mathrm{H}}}
\newcommand{\siemens}{\ensuremath{\mathrm{S}}}

% circuiti
\usepackage{circuitikz}
\usetikzlibrary{babel}

% disegni
\usepackage{pgfplots}
\pgfplotsset{width=10cm,compat=1.9}

% disponi codice
\usepackage{listings}
\usepackage[table]{xcolor}

\lstdefinestyle{codestyle}{
		backgroundcolor=\color{black!5}, 
		commentstyle=\color{codegreen},
		keywordstyle=\bfseries\color{magenta},
		numberstyle=\sffamily\tiny\color{black!60},
		stringstyle=\color{green!50!black},
		basicstyle=\ttfamily\footnotesize,
		breakatwhitespace=false,         
		breaklines=true,                 
		captionpos=b,                    
		keepspaces=true,                 
		numbers=left,                    
		numbersep=5pt,                  
		showspaces=false,                
		showstringspaces=false,
		showtabs=false,                  
		tabsize=2
}

\lstdefinestyle{shellstyle}{
		backgroundcolor=\color{black!5}, 
		basicstyle=\ttfamily\footnotesize\color{black}, 
		commentstyle=\color{black}, 
		keywordstyle=\color{black},
		numberstyle=\color{black!5},
		stringstyle=\color{black}, 
		showspaces=false,
		showstringspaces=false, 
		showtabs=false, 
		tabsize=2, 
		numbers=none, 
		breaklines=true
}

\lstdefinelanguage{javascript}{
	keywords={typeof, new, true, false, catch, function, return, null, catch, switch, var, if, in, while, do, else, case, break},
	keywordstyle=\color{blue}\bfseries,
	ndkeywords={class, export, boolean, throw, implements, import, this},
	ndkeywordstyle=\color{darkgray}\bfseries,
	identifierstyle=\color{black},
	sensitive=false,
	comment=[l]{//},
	morecomment=[s]{/*}{*/},
	commentstyle=\color{purple}\ttfamily,
	stringstyle=\color{red}\ttfamily,
	morestring=[b]',
	morestring=[b]"
}

% disponi sezioni
\usepackage{titlesec}

\titleformat{\section}
	{\sffamily\Large\bfseries} 
	{\thesection}{1em}{} 
\titleformat{\subsection}
	{\sffamily\large\bfseries}   
	{\thesubsection}{1em}{} 
\titleformat{\subsubsection}
	{\sffamily\normalsize\bfseries} 
	{\thesubsubsection}{1em}{}

% disponi alberi
\usepackage{forest}

\forestset{
	rectstyle/.style={
		for tree={rectangle,draw,font=\large\sffamily}
	},
	roundstyle/.style={
		for tree={circle,draw,font=\large}
	}
}

% disponi algoritmi
\usepackage{algorithm}
\usepackage{algorithmic}
\makeatletter
\renewcommand{\ALG@name}{Algoritmo}
\makeatother

% disponi numeri di pagina
\usepackage{fancyhdr}
\fancyhf{} 
\fancyfoot[L]{\sffamily{\thepage}}

\makeatletter
\fancyhead[L]{\raisebox{1ex}[0pt][0pt]{\sffamily{\@title \ \@date}}} 
\fancyhead[R]{\raisebox{1ex}[0pt][0pt]{\sffamily{\@author}}}
\makeatother

\begin{document}
% sezione (data)
\section{Lezione del 07-05-26}

% stili pagina
\thispagestyle{empty}
\pagestyle{fancy}

% testo
Abbiamo iniziato a vedere, nella lezione 19, i \textit{codici a blocco}, e in praticolare i codici a blocco \textbf{lineari}. Ricordiamo brevemente che questi consistevano nel prodotto di \textit{blocchi} di $k$ bit con una certa matrice generatrice $G$, ottenendo \textit{parole di codice} da $n$ bit:
$$
\mathbf{c}_n = \mathbf{b}_k G_{k \times n}
$$

Avevamo quindi dato la definizione di distanza di \textit{Hamming} (definizione 19.1) per dare una norma allo spazio delle parole di codice. Questa ci dava una definizione di \textit{distanza minima} per un codice a blocco, data come la minima distanza di Hamming fra 2 parole di codice (non coincidenti).

\subsection{Codici a blocco lineari sistematici}
Definiamo i codici a blocco lineari \textbf{sistematici} come composti da parole di codice disposte come segue:
\begin{itemize}
	\item Una prima parte, di $k$ bit, contenenti i bit di informazione vera e propria (cioè l'immagine del blocco da $k$ bit);
	\item Una seconda parte data da $n - k$ bit di \textit{parità}. Notiamo che questi non sono chiamati così perché sfruttano necessariamente la codifica in parità dei codici a blocco. Infatti, i contenuti della matrice $P$ possono essere arbitrari.
\end{itemize}

Diciamo che in generale il codice è in forma sistematica quando la matrice generatrice rispetta:
$$
G = [I_k, P]
$$
dove $I_k$ è la matrice identità di dimensioni $k \times k$, e $P$ è la matrice di parità di dimensioni $k \times n - k$.

Si ha quindi che il risultato è disposto nella forma che volevamo:
$$
\mathbf{c} = \mathbf{b} G = \mathbf{b} [I_k, P] = [\mathbf{b}, \mathbf{b} P] = [\mathbf{b}, \mathbf{p}]
$$
cioè dove si hanno $k$ bit del blocco $\mathbf{b}$ originale, e $\mathbf{p}$, cioè un vettore riga da $n - k$ bit di parità.

\par\smallskip

\noindent
\textbf{\textsf{Matrici di controllo}}

\noindent
A partire dalla matrice di parità, possiamo definire la matrice di \textit{controllo} del codice:
$$
H = [P^T, I_{n - k}]
$$
Per ciascuna parola di codice $\mathbf{c}$, si ha che vale:
$$
\mathbf{c} H^T = 0
$$
questo in quanto:
$$
\mathbf{c} H^t = \mathbf{b} G H^T = \mathbf{b}
\begin{pmatrix}
	I_k & P
\end{pmatrix}
\begin{pmatrix}
	{P^T}^T \\ I_{n-k}^T
\end{pmatrix}
= b \left( P + P\right) = 0
$$
in quanto in modulo 2 $P + P = 0$.
Questa matrice risulta quindi utile per verificare la validità di una parola di codice, in quanto se il prodotto $\mathbf{c} H^T$ non fa 0 c'è stato un errore di trasmissione. 

\subsubsection{Codici di Hamming}
I codici di Hamming sono definiti da un parametro $m \geq 2$ che soddisfa:
$$
n = 2^m - 1, \quad k = 2^m - m - 1
$$

A questo punto la matrice di parità $P$ viene costruita così che le colonne della matrice di controllo $H$:
$$
H = [P^T, I_{n - k}]
$$
sia formata in colonne da tutte le possibili $2^m - 1$ combinazioni (esclusa l'$m$-upla di zeri) di $m$ bit.

Si dimostra che la distanza minima Hamming di un qualsiasi codice di Hamming è 3.

\subsubsection{Rivelazione e correzione degli errori nei codici a blocco}
Abbiamo introdotto sistemi di comunicazione basati su:
\begin{itemize}
	\item Source encoder, channel encoder e modulatore a monte (cioè a trasmettitore);
	\item Demodulatore, channel decoder e source decoder a valle, in quest'ordine (cioè ricevitore).
\end{itemize}

Abbiamo che la $n$-upla $\mathbf{y}$ a valle del demodulatore può essere rappresentata come il vettore:
$$
\mathbf{y} = \mathbf{x} + \mathbf{e}
$$
dove $\mathbf{e}$ è il vettore degli errori introdotto dal canale di propagazione (cioè la somma dei rumori e delle distorsioni incontrate), mentre $\mathbf{x}$ era la parola di codice ricavata a monte del channel encoder a trasmettitore.

Supponiamo che il canale introduca un certo numero di errori, tali per cui il peso di Hamming:
$$
w(\mathbf{e}) > 0
$$
è maggiore di zero. La rivelazione degli errori, nel caso si usi un codice a blocco sistematico, si fa semplicemente prendendo la matrice di controllo:
$$
\mathbf{y} H^T \neq 0 \ \implies \ \text{$\mathbf{y}$ non è una parola di codice} 
$$
e si ha un errore \textit{rilevabile}, mentre in caso contrario:
$$
\mathbf{y} H^T = 0 \ \implies \ \text{$\mathbf{y}$ è una parola di codice} 
$$
Notiamo che il fatto che $\mathbf{y}$ sia una parola di codice non rende automaticamente corretto il valore $\mathbf{y}$ ricevuto, cioè non assicura $\mathbf{y} = \mathbf{x}$.
Abbiamo infatti che può accadere il caso dove $\mathbf{y} H^T = 0$ ma comunque $\mathbf{y} \neq \mathbf{x}$. Chiamiamo questo caso quello dell'errore \textit{non rilevabile}.

Possiamo dare il seguente teorema:
\begin{theorem}{Errori rilevabili da codici a blocchi}
	Un codice è in grado di rilevare con certezza fino a $d_{min} - 1$ (dove $d_{min}$ è la distanza di Hamming minima) errori.
\end{theorem}

Questo si dimostra notando che:
\begin{itemize}
	\item Se $d_H(\mathbf{x}, \mathbf{y}) < d_{min}$, allora $\mathbf{y}$ non può essere una parola di codice (l'errore è rilevabile);
	\item Se $d_H(\mathbf{x}, \mathbf{y}) = d_{min}$, allora esiste almeno una parola di codice $\mathbf{c}$ tale che:
$$
d_H(x, c) = d_{min}
$$
e se $\mathbf{y} = \mathbf{c}$, l'errore non è rilevabile. Chiaramente, nel caso $\geq d_{min}$, la situazione è ancora peggiore e gli errori non sono rilevabili. \hfill $\square$
\end{itemize}

\subsubsection{Decodifica a massima verosimiglianza}
La decodifica a massima verosimiglianza (detta anche \textbf{ML}, \textit{Maximum Likelihood}) consiste nel prendere, ricevuto un dato $\mathbf{y}$, la $\mathbf{x}$ che rispetta:
$$
\hat{\mathbf{x}} = \mathrm{argmax}_{x \in \mathcal{C}} P(\mathbf{y} | \mathbf{x})
$$
dove $\mathcal{C}$ è lo spazio di tutte le parole di codice. L'$\mathrm{argmax}$ prende l'argomento (cioè la $\mathbf{x}$ che massimizza), per cui questo consiste nel trovare il vettore $\mathbf{x}$ che fra tutte le $2^k$ possibili parole di codice di lunghezza $n$, massimizza la probabilità condizionata:
$$
P(\mathbf{y} | \mathbf{x})
$$
cioè che si ottenga la parola di codice $\mathbf{x}$ al ricevitore, trasmessa la $\mathbf{y}$.

Cerchiamo un espressione per questa probabilità. Se gli eventi errori (con probabilità $p$) sono indipendenti da bit a bit, si ha che vale:
$$
P(\mathbf{y} | \mathbf{x}) = \prod_{l = 1}^n P(y_l | x_l)
$$
La probabilità $P(y_l | x_l)$ può assumere 2 valori:
$$
P(y_l | x_l) =
\begin{cases}
	1 - p , & \quad P(y_l = x_l | x_l) \\
	p , & \quad P(y_l \neq x_l | x_l)
\end{cases}
$$
dove $p$ è la probabilità di errore sul bit.

Ora, la distanza di Hamming $d_H(\mathbf{x}, \mathbf{y})$ misura il numero di posizioni diverse tra $\mathbf{x}$ e $\mathbf{y}$, per cui $n - d_H(\mathbf{x}, \mathbf{y})$ misura il numero posizioni uguali tra $\mathbf{x}$ e $\mathbf{y}$. Quindi, si ottiene:
$$
P(\mathbf{y} | \mathbf{x}) = p^{d_H(\mathbf{x}, \mathbf{y})} (1 - p)^{n - d_H(\mathbf{x}, \mathbf{y})} = (1 - p)^n \left( \frac{p}{1 - p} \right)^{d_H(\mathbf{x}, \mathbf{y})}
$$
da cui questo diminuisce con $d_H(\mathbf{x}, \mathbf{y})$ (in quanto $\frac{p}{1- p} < 1$ se $p < \frac{1}{2}$), e massimizzare la verosomiglianza $P(\mathbf{y} | \mathbf{x})$ equivale a minimizzare la distanza di Hamming $d_H(\mathbf{x}, \mathbf{y})$.

\end{document}
